\documentclass[11pt]{article}
\usepackage{amsmath, amssymb, geometry, mathtools, booktabs, longtable}
\usepackage{hyperref}
\geometry{margin=1in}

\title{RG Coarse-Eta Flow with Learned Gaussian Priors for 2D $\phi^4$}
\author{Kostas Orginos}
\date{\today}

\begin{document}
\maketitle

\section{Purpose}

This note documents the independent Gaussian-prior evolution branch implemented
in
\texttt{jaxqft/models/phi4\_rg\_coarse\_eta\_gaussian\_flow.py},
with trainer
\texttt{scripts/phi4/train\_rg\_coarse\_eta\_gaussian\_flow.py}
and checker
\texttt{scripts/phi4/check\_rg\_coarse\_eta\_gaussian\_flow.py}.

This branch leaves
\texttt{jaxqft/models/phi4\_rg\_coarse\_eta\_flow.py}
unchanged and adds learned multivariate Gaussian priors on top of the same
coarse-eta Wilsonian RG structure.

\section{High-Level Design}

At each non-terminal RG level:
\begin{enumerate}
\item split the field into a coarse scalar field $c$ and local block
fluctuations $\eta \in \mathbb{R}^3$,
\item recursively trivialize the coarse field,
\item apply a learned Gaussian linear map to the local latent fluctuation
variables,
\item apply the existing nonlinear coarse-eta fluctuation flow,
\item merge the coarse field and fluctuations back to the finer lattice.
\end{enumerate}

The key point is that the coarse field remains fixed during the whole
fluctuation update, so the level map remains Wilsonian in the same sense as the
base coarse-eta model.

\section{Added Flexibility}

The current trainer supports two forms of flexibility that were added after the
initial Gaussian-prior branch was implemented.

\subsection{Per-Level Nonterminal Capacity}

For a lattice of linear size $L=2^K$, there are $K-1$ non-terminal RG levels.
They are ordered in the code from finest to coarsest. For example:
\begin{itemize}
\item for $L=16$, the non-terminal coarse lattices are $8^2$, $4^2$, $2^2$,
\item for $L=32$, the non-terminal coarse lattices are $16^2$, $8^2$, $4^2$,
$2^2$.
\end{itemize}

The trainer can now assign different capacities to these levels:
\begin{itemize}
\item \texttt{width\_levels}: nonlinear eta-flow conditioner widths,
\item \texttt{n\_cycles\_levels}: red/black cycle counts,
\item \texttt{radius\_levels}: nonlinear conditioner patch radii,
\item \texttt{gaussian\_width\_levels}: widths of the conditioned Gaussian MLP,
\item \texttt{gaussian\_radius\_levels}: Gaussian-conditioning patch radii.
\end{itemize}

If these lists are not provided, the code falls back to the uniform scalar
values \texttt{width}, \texttt{n\_cycles}, \texttt{radius},
\texttt{gaussian\_width}, and \texttt{gaussian\_radius}.

\subsection{TOML-Driven Runs and Staged Training}

The trainer supports \texttt{--config <file.toml>}. A TOML file can define:
\begin{itemize}
\item fresh runs or resume runs,
\item physics parameters,
\item model architecture,
\item validation settings,
\item automated batch-ramp and learning-rate schedules.
\end{itemize}

The supported schedule forms are:
\begin{itemize}
\item explicit stages:
\begin{verbatim}
[[schedule.stage]]
epoch_end = 1000
batch = 32
lr = 3e-4
\end{verbatim}
\item generated batch-doubling ramps:
\begin{verbatim}
[schedule.ramp]
initial_batch = 32
epochs_per_stage = 1000
num_doubles = 5
lr = 3e-4
\end{verbatim}
\item late-stage learning-rate anneals:
\begin{verbatim}
[schedule.anneal]
epoch_ends = [7000, 8000]
lrs = [1e-4, 3e-5]
batch = 1024
\end{verbatim}
\end{itemize}

When a schedule is present, the trainer derives the total target epoch from the
last schedule stage, switches batch size and learning rate automatically at
stage boundaries, and preserves the schedule inside the checkpoint payload.

Validation can also run at every stage boundary:
\begin{verbatim}
[validation]
enabled = true
each_stage = true
batch = 256
super = 4
\end{verbatim}
This stage validation uses a folded random key derived from the training RNG,
so enabling it does not perturb the training trajectory.

\section{RG Split}

Let the level-$k$ lattice be $\Lambda_k$ with linear size $L_k = L / 2^k$.
Each $2 \times 2$ block is written as
\[
v = (v_{00}, v_{01}, v_{10}, v_{11}) \in \mathbb{R}^4.
\]

For \texttt{rg\_type="average"} the code uses the orthonormal Haar basis
\[
(c,\eta_1,\eta_2,\eta_3) = H v,
\]
with
\[
H=
\begin{pmatrix}
\tfrac12 & \tfrac12 & \tfrac12 & \tfrac12 \\
\tfrac12 & -\tfrac12 & \tfrac12 & -\tfrac12 \\
\tfrac12 & \tfrac12 & -\tfrac12 & -\tfrac12 \\
\tfrac12 & -\tfrac12 & -\tfrac12 & \tfrac12
\end{pmatrix}.
\]
For \texttt{rg\_type="select"} the code uses
\[
c=v_{00}, \qquad \eta=(v_{01},v_{10},v_{11}).
\]

Thus the level state should be viewed as a field on the coarse lattice
\[
y_b = (c_b,\eta_b), \qquad c_b \in \mathbb{R}, \quad \eta_b \in \mathbb{R}^3,
\qquad b \in \Lambda_{k+1}.
\]

\section{Latent Prior and Gaussian Maps}

The global latent prior remains sitewise standard normal on the original
lattice:
\[
z \sim \mathcal{N}(0, I).
\]

The learned Gaussian priors are implemented as exact lower-triangular linear
maps acting on standard-normal latent variables.

\subsection{Intermediate $3 \times 3$ Gaussian}

At each non-terminal RG level, before the nonlinear eta flow, the latent
fluctuation vector $z_{\eta,b} \in \mathbb{R}^3$ is mapped to an intermediate
variable
\[
\tilde z_{\eta,b} = L_b(c)\, z_{\eta,b},
\]
where $L_b(c)$ is lower triangular with positive diagonal. In the code this is
parameterized by diagonal log-scales and three strictly lower-triangular
entries:
\[
L_b =
\begin{pmatrix}
e^{\alpha_1} & 0 & 0 \\
\ell_{21} & e^{\alpha_2} & 0 \\
\ell_{31} & \ell_{32} & e^{\alpha_3}
\end{pmatrix}.
\]

There are three modes:
\begin{itemize}
\item \texttt{eta\_gaussian = none}: $L_b = I$.
\item \texttt{eta\_gaussian = level}: one shared learned $L^{(\ell)}$ per RG
level.
\item \texttt{eta\_gaussian = coarse\_patch}: $L_b$ is conditioned on a local
patch of coarse fields only.
\end{itemize}

For the conditioned mode, the context is
\[
u_b = \tilde c|_{\mathcal N_{r_g}(b)},
\]
where $\tilde c_b$ is the physical coarse scalar used for conditioning and
$\mathcal N_{r_g}(b)$ is the coarse-lattice patch of radius
\texttt{gaussian\_radius}. The conditioner is a local patch MLP. It depends
only on coarse fields, not on neighboring $\eta$ variables.

\subsection{Terminal $4 \times 4$ Gaussian}

At the final $2 \times 2$ block, the code uses a terminal RealNVP and an
optional learned zero-mean $4$-variate Gaussian. The generative map is
\[
z_{\mathrm{term}} \xmapsto{L_{\mathrm{term}}}
\tilde z_{\mathrm{term}} \xmapsto{\mathrm{RealNVP}}
x_{\mathrm{term}}.
\]

Here
\[
L_{\mathrm{term}} \in \mathbb{R}^{4 \times 4}
\]
is lower triangular with positive diagonal. The mode is selected by
\texttt{terminal\_prior}:
\begin{itemize}
\item \texttt{std}: $L_{\mathrm{term}} = I$,
\item \texttt{learned}: full learned lower-triangular $4 \times 4$ map.
\end{itemize}

\section{Nonlinear Coarse-Eta Flow}

After the Gaussian map, the code applies the same nonlinear coarse-eta flow as
the base model. Suppressing the level index, the coarse lattice is split into
red and black checkerboards:
\[
\Lambda = R \cup B.
\]

For one red sweep:
\begin{itemize}
\item all coarse scalars $c$ are visible and fixed,
\item black fluctuations are visible and frozen,
\item red fluctuations are updated.
\end{itemize}

At a red site $b$, the update is a conditional lower-triangular affine map on
the local three-component fluctuation vector:
\[
\eta_b' = A_b(h)\, \eta_b + t_b(h),
\]
where $A_b(h)$ is lower triangular with positive diagonal, and the conditioner
feature $h$ is produced by a local patch MLP from
\[
u_b = \bigl(\eta^{\mathrm{masked}}, \tilde c \bigr)\big|_{\mathcal N_r(b)}.
\]
The black sweep is analogous. Repeating red then black gives one cycle, and the
number of cycles is controlled by \texttt{n\_cycles}.

\section{Full Generative and Inverse Maps}

Let $G_\ell^{\mathrm{coarse}}$ denote the recursive coarse-field map at level
$\ell$, $G_\ell^{\mathrm{gauss}}$ the Gaussian linear map on $\eta$, and
$G_\ell^{\eta}$ the nonlinear coarse-eta flow. Then for a non-terminal level,
\[
G_\ell(z_\ell) =
\operatorname{merge}\!\Bigl(
G_{\ell+1}^{\mathrm{coarse}}(z_{\ell+1}),
G_\ell^{\eta}\bigl(
G_\ell^{\mathrm{gauss}}(z_{\eta,\ell}) \,\big|\, x_{\ell+1}
\bigr)
\Bigr),
\]
with
\[
(z_{\ell+1}, z_{\eta,\ell}) = \operatorname{split}(z_\ell),
\qquad
x_{\ell+1} = G_{\ell+1}^{\mathrm{coarse}}(z_{\ell+1}).
\]

At the terminal level,
\[
G_{\mathrm{term}} = \mathrm{RealNVP} \circ L_{\mathrm{term}}.
\]

The inverse order is reversed:
\[
F_\ell = \operatorname{split}^{-1} \circ
\Bigl(
F_{\ell+1}^{\mathrm{coarse}},
\left(G_\ell^{\mathrm{gauss}}\right)^{-1}
\circ
\left(G_\ell^{\eta}\right)^{-1}
\Bigr).
\]

\section{Exact Log Density}

Because every piece is triangular or a standard coupling flow, the log-density
is exact. If
\[
z = F(x), \qquad \log p_Z(z) = -\frac12 \|z\|^2 - \frac{N}{2}\log(2\pi),
\]
then
\[
\log p_X(x) = \log p_Z(z) + \log \left| \det \frac{\partial F}{\partial x} \right|.
\]

Each intermediate $3 \times 3$ Gaussian contributes
\[
\log \left| \det \frac{\partial F_{\mathrm{gauss}}}{\partial x} \right|
= -(\alpha_1 + \alpha_2 + \alpha_3),
\]
and the terminal $4 \times 4$ Gaussian contributes
\[
\log \left| \det \frac{\partial F_{\mathrm{term,gauss}}}{\partial x} \right|
= -(\alpha_1 + \alpha_2 + \alpha_3 + \alpha_4).
\]

\section{Parity Symmetry}

The branch supports
\[
\phi \mapsto -\phi, \qquad z \mapsto -z
\]
through \texttt{parity = sym}. For the conditioned intermediate Gaussian, this
is enforced by symmetrizing the raw conditioner outputs as an even function of
the coarse-field context:
\[
\mathrm{raw}(u) \mapsto \frac12 \bigl(\mathrm{raw}(u) + \mathrm{raw}(-u)\bigr).
\]
Since the Gaussian is zero mean, only even dependence is needed there.

\section{Trainer CLI}

The trainer is
\texttt{scripts/phi4/train\_rg\_coarse\_eta\_gaussian\_flow.py}.

\subsection*{Physics and geometry}
\begin{longtable}{@{}ll@{}}
\toprule
Flag & Meaning \\
\midrule
\texttt{--L} & lattice size, requires a square power of two \\
\texttt{--mass} & bare quadratic coupling \\
\texttt{--lam} & quartic coupling \\
\texttt{--rg-type \{average,select\}} & block decomposition rule \\
\bottomrule
\end{longtable}

\subsection*{Flow architecture}
\begin{longtable}{@{}ll@{}}
\toprule
Flag & Meaning \\
\midrule
\texttt{--width} & hidden width of the nonlinear coarse-eta conditioner \\
\texttt{--width-levels} & comma-separated per-level widths, finest $\to$ coarsest \\
\texttt{--n-cycles} & number of red/black coarse-lattice eta cycles per level \\
\texttt{--n-cycles-levels} & comma-separated per-level cycle counts \\
\texttt{--radius} & nonlinear conditioner patch radius on the coarse lattice \\
\texttt{--radius-levels} & comma-separated per-level nonlinear radii \\
\texttt{--eta-gaussian \{none,level,coarse\_patch\}} & intermediate Gaussian mode \\
\texttt{--gaussian-radius} & conditioning radius for \texttt{coarse\_patch} Gaussian \\
\texttt{--gaussian-radius-levels} & comma-separated per-level Gaussian radii \\
\texttt{--gaussian-width} & hidden width for the conditioned Gaussian MLP \\
\texttt{--gaussian-width-levels} & comma-separated per-level Gaussian widths \\
\texttt{--terminal-prior \{std,learned\}} & terminal $4$-variate Gaussian mode \\
\texttt{--terminal-n-layers} & terminal RealNVP depth \\
\texttt{--terminal-width} & terminal RealNVP width \\
\texttt{--output-init-scale} & small output initialization scale \\
\texttt{--parity \{none,sym\}} & optional $Z_2$ symmetry projection \\
\texttt{--log-scale-clip} & clip on diagonal log-scales \\
\texttt{--offdiag-clip} & clip on lower-triangular off-diagonal entries \\
\bottomrule
\end{longtable}

\subsection*{Training and checkpointing}
\begin{longtable}{@{}ll@{}}
\toprule
Flag & Meaning \\
\midrule
\texttt{--config} & TOML input card for fresh or resume runs \\
\texttt{--epochs} & target epoch count; on resume this is the absolute target \\
\texttt{--batch} & batch size; preserved from checkpoint unless overridden \\
\texttt{--lr} & learning rate; preserved from checkpoint unless overridden \\
\texttt{--resume} & checkpoint to load \\
\texttt{--save} & checkpoint output file \\
\texttt{--save-every} & periodic checkpoint cadence \\
\texttt{--validate} & run validation at the end of the script \\
\texttt{--validate-each-stage} & when using a schedule, validate at each stage boundary \\
\texttt{--validate-batch} & validation mini-batch size \\
\texttt{--validate-super} & number of validation batches to aggregate \\
\bottomrule
\end{longtable}

\section{Input Files}

The repository currently provides three useful TOML cards for this branch:
\begin{itemize}
\item \texttt{configs/phi4/rg\_coarse\_eta\_gaussian\_fresh.toml}
\begin{itemize}
\item example fresh run,
\item uses the canonical $16^2$ physical point,
\item uses the batch ramp $32 \to 64 \to 128 \to 256 \to 512 \to 1024$ and
then a two-stage learning-rate anneal.
\end{itemize}
\item \texttt{configs/phi4/rg\_coarse\_eta\_gaussian\_resume.toml}
\begin{itemize}
\item example resume card,
\item intended for continuing a saved checkpoint while preserving the same
model family and staged training logic.
\end{itemize}
\item \texttt{configs/phi4/rg\_coarse\_eta\_gaussian\_L32\_perlevel.toml}
\begin{itemize}
\item direct $32^2$ extension of the successful per-level $16^2$ setup,
\item keeps the simple learned terminal Gaussian and terminal RealNVP,
\item uses per-level settings
\[
\texttt{width\_levels} = [128,96,64,96], \quad
\texttt{n\_cycles\_levels} = [3,3,2,3], \quad
\texttt{radius\_levels} = [1,1,1,2].
\]
\end{itemize}
\end{itemize}

The fresh/resume split is controlled by
\[
\texttt{mode} \in \{\texttt{"fresh"},\texttt{"resume"}\}.
\]
In resume mode the file must provide
\[
\texttt{[io] resume = "..."}.
\]
Regular CLI flags still override TOML defaults.

\section{Checker CLI}

The checker is
\texttt{scripts/phi4/check\_rg\_coarse\_eta\_gaussian\_flow.py}. It provides:
\begin{itemize}
\item invertibility test,
\item Jacobian test against autodiff,
\item timing breakdown.
\end{itemize}

Main checker-only knobs:
\begin{itemize}
\item \texttt{--shape H,W}
\item \texttt{--jac-shape H,W}
\item \texttt{--tests invertibility,jacobian,timing}
\item \texttt{--invert-tol}
\item \texttt{--jac-tol}
\item \texttt{--jac-perturb-scale}
\item \texttt{--timing-warmup}
\item \texttt{--timing-iters}
\item \texttt{--selfcheck-fail}
\end{itemize}

\section{Example Commands}

Fresh run with the cheaper shared-per-level Gaussian:
\begin{verbatim}
source /opt/python/jax/bin/activate
MPLCONFIGDIR=/tmp/mpl-cache JAX_PLATFORMS=cpu \
python scripts/phi4/train_rg_coarse_eta_gaussian_flow.py \
  --L 16 --mass -0.4 --lam 2.4 \
  --width 64 --n-cycles 2 --radius 1 \
  --eta-gaussian level \
  --terminal-prior learned \
  --batch 32 --lr 3e-4 --parity sym \
  --epochs 1000 --validate
\end{verbatim}

Fresh run with the coarse-conditioned Gaussian:
\begin{verbatim}
source /opt/python/jax/bin/activate
MPLCONFIGDIR=/tmp/mpl-cache JAX_PLATFORMS=cpu \
python scripts/phi4/train_rg_coarse_eta_gaussian_flow.py \
  --L 16 --mass -0.4 --lam 2.4 \
  --width 64 --n-cycles 2 --radius 1 \
  --eta-gaussian coarse_patch \
  --gaussian-radius 1 --gaussian-width 64 \
  --terminal-prior learned \
  --batch 32 --lr 3e-4 --parity sym \
  --epochs 1000 --validate
\end{verbatim}

Checker:
\begin{verbatim}
source /opt/python/jax/bin/activate
MPLCONFIGDIR=/tmp/mpl-cache JAX_PLATFORMS=cpu \
python scripts/phi4/check_rg_coarse_eta_gaussian_flow.py --selfcheck-fail
\end{verbatim}

Fresh TOML-driven run:
\begin{verbatim}
source /opt/python/jax/bin/activate
MPLCONFIGDIR=/tmp/mpl-cache JAX_PLATFORMS=cpu \
python scripts/phi4/train_rg_coarse_eta_gaussian_flow.py \
  --config configs/phi4/rg_coarse_eta_gaussian_fresh.toml
\end{verbatim}

Resume a staged run from checkpoint:
\begin{verbatim}
source /opt/python/jax/bin/activate
MPLCONFIGDIR=/tmp/mpl-cache JAX_PLATFORMS=cpu \
python scripts/phi4/train_rg_coarse_eta_gaussian_flow.py \
  --config configs/phi4/rg_coarse_eta_gaussian_L32_perlevel.toml \
  --resume rg_coarse_eta_gauss_L32_m-0.4_l2.4_perlevel_w128-96-64-96_\
nc3-3-2-3_r1-1-1-2_eglevel_tglearned.pkl
\end{verbatim}

\section{Canonical $16^2$ Comparison}

At the canonical point
\[
L=16, \qquad m^2=-0.4, \qquad \lambda=2.4,
\]
both intermediate Gaussian variants were run with the batch ramp
\[
32 \to 64 \to 128 \to 256 \to 512 \to 1024
\]
through epoch $6000$, using \texttt{terminal\_prior=learned} in both cases.

\begin{center}
\begin{tabular}{lccc}
\toprule
Intermediate Gaussian & $\mathrm{std}(\Delta S)$ & $\mathrm{std}(w)$ & ESS \\
\midrule
\texttt{level} & 0.8416 & 0.9698 & 0.5153 \\
\texttt{coarse\_patch} & 0.8619 & 0.9871 & 0.5065 \\
\bottomrule
\end{tabular}
\end{center}

So in the current implementation the coarse-conditioned intermediate Gaussian is
not a clear improvement over the cheaper shared-per-level Gaussian.

\section{Current Interpretation}

The Gaussian-prior branch is useful because it cleanly separates:
\begin{itemize}
\item a standard-normal global latent prior,
\item learned local Gaussian whitening at each RG level,
\item the nonlinear fluctuation trivialization map.
\end{itemize}

The present tests suggest:
\begin{itemize}
\item a learned terminal $4$-variate Gaussian is worthwhile,
\item a learned intermediate $3 \times 3$ Gaussian also works well,
\item conditioning that intermediate Gaussian only on coarse fields does not
yet produce a clear gain over a simpler levelwise covariance.
\end{itemize}

\end{document}
